\scanpage{40}\mspage{27}

$3^\circ$) \inserted{Supposons encore $I$ artinien\textsuperscript{*}.} Soit $W\unsure{\subset|\in}\mathrm{Fl}(M)$ un localiseur modérément \deleted{\unsure{satisfait}}\unsure{saturé|saturisé}, tel que \deleted{\unsure{la}} $TC\overset{\mathrm{def}}{=}C\cap W$ satisfasse également aux conditions de stabilité (1.1.) (par \unsure{cochgt} de base, par composition transfinie).

% slanted
\margin{J'ignore si cette hypothèse est nécessaire, cf commentaire p. \unsure{39|35}.}

Les conditions suivantes sur une flèche $u_\bullet:\deleted{\unsure{Y}}X_\bullet\longrightarrow\unsure{Y|Z}$ dans $M(I^{\unsure{o|\circ}})$ sont équivalentes

a) $u_\bullet\in(TC)_I$, i.e. pour tout $i_0\in I$ et $K$ ouvert de $J=I_{/i_0}$, la flèche $u_{K,J}$ est dans $TC$

b) $u_\bullet\in C_I\cap W_{I^o}$, i.e. \deleted{Pour les} pour tout $i_0\in I$, on a $u_{i_0}\in W$, et pour tout ouvert $K$ de $I_{/i_0}=J$, on a $u_{K,J}\in C$.

Donnons tout de suite les corollaires principaux :

\uline{Corollaire \unsure{2.2.|2.1.}}: Supposons $(u_\bullet:X_\bullet\longrightarrow Y)\in C_I$. Alors

a) $\forall\,i\in I$, $u_i:X_i\longrightarrow Y_i$ est dans $C$.

b) \deleted{\unsure{$\forall$}} La flèche

\[
 u_I:X_I\longrightarrow Y_I\tag*{(2.2.1.)}
\]

est dans $C$. Plus généralement, pour \unsure{tte} partie ouverte $J$ de $I$, la flèche

\[
 u_J:X_J\longrightarrow Y_J\tag*{(2.2.2)}
\]

est dans $C$.

% archive sheet number: 39

